\documentclass[12pt,a4paper]{letter}
\usepackage[margin=2.5cm]{geometry}
\usepackage{hyperref}
\usepackage{microtype}

\signature{Nisong Monyimba, Vincent Pizziconi, Aurel Coza}
\address{School of Biological and Health Systems Engineering\\
Arizona State University, Tempe, AZ, USA\\
nisongmonyimba278@gmail.com}
\date{2026}

\begin{document}

\begin{letter}{
    The Editorial Board\\
    Engineering with Computers\\
    Springer Nature
}

\opening{Dear Editors,}

We submit for your consideration the manuscript:

\begin{center}
\textbf{\texttt{brinkgrad}: Open-source topology optimisation for coupled flow-transport in FEniCSx}\\[4pt]
{\normalfont\textit{A software implementation for Brinkman--convection-diffusion systems}}
\end{center}

\textbf{Scope and fit.}
This manuscript falls under the \emph{Engineering Software Development}
category of Engineering with Computers.
\texttt{brinkgrad} is an openly distributed, CI-verified, and
Docker-reproducible FEniCSx framework for topology optimisation of
coupled Brinkman--convection-diffusion systems targeting
concentration-gradient generation in microfluidic channels.
We are not aware of any prior open-source FEniCSx package providing
this capability with SUPG-stabilised adjoint transport, integrated CI,
and single-command reproducibility.

\textbf{Principal contributions.}
\begin{enumerate}
\item \textbf{SUPG-stabilised coupled adjoint framework.}
At the P\'{e}clet numbers relevant to microfluidic transport
($\mathrm{Pe}\sim10^2$--$10^5$), both the forward and adjoint
convection-diffusion equations require stabilisation.
\texttt{brinkgrad} applies SUPG consistently to both equations,
with gradient accuracy verified by directional derivative,
descent, and Taylor remainder tests.

\item \textbf{Coupled Brinkman--convection-diffusion topology optimisation.}
The framework targets concentration-mismatch objectives for
coupled flow-transport systems --- extending the standard
Borrvall--Petersson~(2003) Stokes-drag framework to include
species transport.
This class of problem is underexplored in the open-source
topology optimisation literature.

\item \textbf{End-to-end reproducibility infrastructure.}
All results are recoverable from a single Docker command
(\texttt{make reproduce}) with a pinned \texttt{dolfinx:v0.7.3} image,
18 automated tests on every commit, and permanent Zenodo archival
(DOI: \href{https://doi.org/10.5281/zenodo.20538833}{10.5281/zenodo.20538833}).

\item \textbf{Framework generality.}
Three target profiles (linear, step, double-peak) and a
volume-fraction sweep demonstrate that the framework adapts
to arbitrary concentration targets without modifying the
adjoint or optimiser structure.
\end{enumerate}

\textbf{Validation.}
Gradient correctness is established through a directional derivative test
and a Taylor remainder test (slope~$\approx1$, consistent with the
$L^2$ Riesz representative formulation).
A four-mesh convergence study (Table~4 in manuscript) and heuristic
baseline comparison (81\% RMSE improvement over the best hand-designed
geometry) demonstrate that the adjoint-based optimiser adds genuine value.

\textbf{Suggested reviewers.}
\begin{itemize}
\item Prof.\ Ole Sigmund (DTU) --- topology optimisation methodology
\item Prof.\ Boyan Lazarov (DTU/LLNL) --- Helmholtz filtering, FEniCS
\item Prof.\ Kurt Maute (CU Boulder) --- topology optimisation, adjoint methods
\item Dr.\ Florian Wechsung (NYU) --- FEniCSx adjoint methods
\end{itemize}

\textbf{Declarations.}
This manuscript has not been published elsewhere and is not under
consideration by any other journal.
All authors have approved the submission.
The code is publicly available at
\url{https://github.com/NisongMonyimba/brinkgrad}
(Zenodo DOI: 10.5281/zenodo.20538833).

\closing{Yours sincerely,}

\end{letter}
\end{document}
